\documentclass[11pt]{article}
\usepackage[a4paper,margin=1in]{geometry}
\usepackage{amsmath,amssymb}
\usepackage{hyperref}
\title{Discrete ODF and Pole-Figure Foundations in PyTex}
\author{PyTex Contributors}
\date{\today}

\begin{document}
\maketitle

\section{Scope}

The current PyTex texture layer implements the foundational algorithms needed before full inversion and harmonic reconstruction:

\begin{itemize}
  \item pole-figure synthesis from orientation sets,
  \item inverse-pole-figure synthesis from orientation sets,
  \item kernel-weighted ODF evaluation,
  \item volume-fraction queries around a reference orientation.
\end{itemize}

\section{Pole-Figure Construction}

Given a pole normal in the crystal frame and an orientation set $\{g_i\}$, the corresponding specimen directions are obtained by applying each orientation to the pole family. The current implementation uses explicit crystal-symmetry expansion of the pole family and stores the resulting specimen directions together with their weights.

\section{ODF Evaluation}

PyTex currently treats the ODF as a weighted orientation set together with an angular kernel. This is a foundational representation rather than a final inversion framework. Two kernels are presently supported:

\begin{itemize}
  \item a halfwidth-calibrated cosine-power kernel exposed as \texttt{de\_la\_vallee\_poussin},
  \item a von Mises--Fisher style kernel.
\end{itemize}

\section{Current Limits}

Full harmonic ODF expansion, experimental PF inversion, and rigorous kernel normalization on $SO(3)$ remain future work. The current representation is deliberately explicit and testable so later algorithms can be layered on top of a stable semantic core.

\section{Normative References}

\begin{thebibliography}{9}
\bibitem{bunge}
H.-J. Bunge,
\textit{Texture Analysis in Materials Science: Mathematical Methods},
Butterworths, 1969.
DOI: \url{https://doi.org/10.1016/C2013-0-11769-2}.

\bibitem{fisher}
R. A. Fisher,
``Dispersion on a Sphere'',
\textit{Proceedings of the Royal Society A}, 217(1130), 1953, 295--305.
DOI: \url{https://doi.org/10.1098/rspa.1953.0064}.
\end{thebibliography}

\end{document}
